\documentclass[12pt, a4paper]{article}

% Pakiety podstawowe
\usepackage[utf8]{inputenc}
\usepackage[T1]{fontenc}
\usepackage[polish]{babel}
\usepackage{lmodern}

% Pakiety do grafiki i tabel
\usepackage{graphicx} % Do wstawiania wykresów
\usepackage{float}    % Do precyzyjnego pozycjonowania (opcja [H])
\usepackage{booktabs} % Do ładniejszych tabel
\usepackage{geometry} % Ustawienia marginesów
\geometry{a4paper, margin=2.5cm}

% Pakiety matematyczne
\usepackage{amsmath}

% Pakiet do linków
\usepackage{hyperref}
\hypersetup{
    colorlinks=true,
    linkcolor=blue,
    filecolor=magenta,      
    urlcolor=cyan,
}

% --- METADANE DOKUMENTU ---
\title{Projektowanie i analiza algorytmów\\ \large Raport: Algorytmy sortowania}
\author{Rafał Jakimowicz \\ 284082}
\date{Marzec 2026}

\begin{document}

\maketitle
\tableofcontents
\newpage

% --- 1. WPROWADZENIE ---
\section{Wprowadzenie}
W niniejszym raporcie przedstawiono implementację oraz analizę efektywności trzech wybranych algorytmów sortowania: sortowanie przez scalanie, Quicksort, sortowanie introspektywne. Celem było zbadanie czasu działania algorytmów w zależności od rozmiaru danych wejściowych oraz ich początkowego stopnia uporządkowania. Algorytmy zostały zaimplementowane w języku C++. Uporządkowanie danych oraz stworzenie wykresów zostało wykonane w Python. Cały kod źródłowy został opublikowany na publiczne  \href{https://github.com/RafalJakimowicz/PiAASortowanie}{repozytorium Github}

% --- 2. OPIS ALGORYTMÓW ---
\section{Opis badanych algorytmów}
Poniżej zostanie opisana teoretyczna charakterystyka wybranych metod sortowania, oraz opisana ich teoretyczna złożoność czasowa algorytmów dla przypadku średniego i najgorszego.

\subsection{Algorytm 1: Sortowanie przez scalanie}
Algorytm działający zgodnie z zasadą "Devide and Conquer", dzielący sortowaną strukturę na coraz mniesze fragmenty za każdym razem na połowy z wykorzystaniem rekurencji które są poźniej scalane w posortowaną strukturę. 
\begin{itemize}
    \item Złożoność czasowa (przypadek średni): $\mathcal{O}(N\log{N})$
    \item Złożoność czasowa (przypadek najgorszy): $\mathcal{O}(N\log{N})$
\end{itemize}

\subsection{Algorytm 2: Quicksort}
Algorytm działający zgodnie z zasadą "Devide and Conquer", w strukturze wybierany jest element i po jego lewej stronie są umieszczane elementy nie większe od niego a po prawej elementy większe, co powoduje poprawne ułożenie danego elementu w tablicy, element ten nosi nazwę pivot. Algorytm jest powtarzany rekurencyjnie dla liczb wiekszych i mniejszych od pivot. Implemenatcja wybiera pivot metodą "Mediany z trzech".
\begin{itemize}
    \item Złożoność czasowa (przypadek średni): $\mathcal{O}(N\log{N})$
    \item Złożoność czasowa (przypadek najgorszy): $\mathcal{O}(n^2)$
\end{itemize}

\subsection{Algorytm 3: sortowanie introspektywne}
Algorytm hybrydowy bazujący na algorytmie Quicksort, implemetujący również sortowanie przez wstawianie i sortowanie przez kopcowanie. Algorytm mierzy warstwy rekurencji sortowania Quicksort i jeśli głębokość przekroczy wartość optymalną zmienia sie na sortowanie przez kopcowanie dla tego fragmentu danych. Dodatkową właściwością jest również zamiana algorytmu sortowania na sortowanie przez wstawianie dla fragmentów poniżej danej liczby elementów. 
\begin{itemize}
    \item Złożoność czasowa (przypadek średni): $\mathcal{O}(N\log{N})$
    \item Złożoność czasowa (przypadek najgorszy): $\mathcal{O}(N\log{N})$
\end{itemize}

% --- 3. EKSPERYMENTY I WYNIKI ---
\section{Przebieg eksperymentów i wyniki}
% Zgodnie z wytycznymi: omówienie przebiegu eksperymentów i wyniki w postaci tabel/wykresów.
Eksperymenty przeprowadzono dla 100 losowych tablic o rozmiarach: 10 000, 50 000, 100 000, 500 000 oraz 1 000 000. Badano czas wykonania dla tablic losowych, częściowo posortowanych (od 25\% do 99,7\%) oraz posortowanych w odwrotnej kolejności.

\subsection{Wyniki dla tablic całkowicie losowych}
Tutaj opisz, jak zachowały się algorytmy. Poniżej przykład wstawienia wykresu wygenerowanego w Pythonie:

% Przykład wstawienia obrazka (odkomentuj i zmień nazwę pliku, gdy wygenerujesz wykres)
% \begin{figure}[H]
%     \centering
%     \includegraphics[width=0.8\textwidth]{wykres_losowe.png}
%     \caption{Porównanie czasów działania dla danych losowych}
%     \label{fig:losowe}
% \end{figure}

\subsection{Wpływ początkowego uporządkowania danych}
Tutaj wstaw i skomentuj wykresy dla różnego procentu posortowania oraz dla tablic ułożonych odwrotnie.

% --- 4. PODSUMOWANIE I WNIOSKI ---
\section{Podsumowanie i wnioski}
% Zgodnie z wytycznymi: wnioski i wyjaśnienie niezgodności z przewidywaniami.
Podsumowanie uzyskanych wyników. 
W tej sekcji odnieś się do tego, czy czasy zmierzone w eksperymentach pokrywają się z teoretyczną złożonością algorytmów opisaną w rozdziale 2. Jeśli np. QuickSort zadziałał gorzej na posortowanych danych (osiągając czas $\mathcal{O}(N^2)$), należy to tutaj wyjaśnić.

% --- 5. LITERATURA ---
\section{Literatura}
% Zgodnie z wytycznymi: literatura i materiały.
\begin{thebibliography}{9}
    \bibitem{jelen} Jeleń Ł., \textit{Notatki do wykładu}. %
    \bibitem{cormen} Cormen T., Leiserson C.E., Rivest R.L., Stein C., \textit{Wprowadzenie do algorytmów}, Wydawnictwo Naukowe PWN. %
    \bibitem Dokumentacja C++.
\end{thebibliography}

\end{document}